\documentclass[11pt,letterpaper]{article}
% Common article worksheet layout. Load after \documentclass.
\usepackage[margin=0.25in,headheight=0pt,headsep=0pt,footskip=0pt]{geometry}
\usepackage{amsmath,amssymb}
\usepackage{tikz,graphicx}
\definecolor{mtuGreen}{RGB}{0,79,42}
\definecolor{mtuGreenLight}{RGB}{60,133,90}
\definecolor{mtuGreenPale}{RGB}{223,236,227}
\newcommand{\vect}[1]{\mathbf{#1}}
\setlength{\parindent}{0pt}
\setlength{\parskip}{8pt}
\pagestyle{empty}
% Fixed name-box dimensions and page position shared by every worksheet.
\newcommand{\WorksheetNameBox}{%
  \begingroup\setlength{\fboxsep}{3pt}\setlength{\fboxrule}{0.4pt}%
  \fbox{\rule[-0.30in]{0pt}{0.42in}\makebox[4.24in][l]{Name:}}%
  \endgroup}
\newcommand{\WorksheetHeading}[1]{%
  \noindent{\small\textbf{MA3521 --- #1}}\hfill\WorksheetNameBox\par}
% A blank entry with fixed writing room, usable inside a matrix.
\newcommand{\AnswerEntry}{\rule[-0.17in]{0pt}{0.47in}\makebox[0.67in]{}}
\newcommand{\AnswerVector}{\begin{pmatrix}\AnswerEntry\\\AnswerEntry\\\AnswerEntry\end{pmatrix}}
\newcommand{\AnswerMatrix}{\begin{pmatrix}
  \AnswerEntry&\AnswerEntry&\AnswerEntry\\
  \AnswerEntry&\AnswerEntry&\AnswerEntry\\
  \AnswerEntry&\AnswerEntry&\AnswerEntry
\end{pmatrix}}

\setlength{\parskip}{5pt}

% Box sizes below are the writing area; the frame adds about 0.1in.
% Change these values to resize the matching boxes throughout the worksheet.
\newcommand{\MatrixSize}{1.3in}
\newcommand{\VectorWidth}{0.8in}
% Matrix writing region is square; these widths leave room for parentheses.
\newcommand{\MatrixWritingSize}{1.05in}
\newcommand{\MatrixBoxWidth}{1.5in}
\newcommand{\VectorWritingWidth}{0.35in}
\newcommand{\AugmentedBoxWidth}{2in}
\newcommand{\AugmentedColumnWidth}{0.35in}
\newcommand{\EigenvalueWidth}{1in}
\newcommand{\ShortHeight}{0.35in}
\newcommand{\SolutionWidth}{1.75in}
\newcommand{\SolutionHeight}{0.85in}
\newcommand{\DerivativeWidth}{2in}
\newcommand{\WideDerivativeWidth}{3in}
\newcommand{\VectorSolutionWidth}{1in}
\newcommand{\LinearSolutionWidth}{6.3in}
\newcommand{\LinearSolutionHeight}{1.4in}
\newcommand{\ParticularSolutionWidth}{4.8in}
\newcommand{\GeneralSolutionWidth}{5.2in}
\newcommand{\LongSolutionHeight}{0.8in}
\newcommand{\SolutionGap}{0.12in}

% One standard box: width, height, contents. Use the named boxes below in math.
\newcommand{\AnswerBox}[3]{%
  \vcenter{\hbox{\fbox{\parbox[c][#2][c]{#1}{\centering #3}}}}}
\newcommand{\BlankBox}[2]{\AnswerBox{#1}{#2}{\strut}}
% Center the invisible writing region on the math axis before sizing delimiters.
\newcommand{\WritingArea}[1]{%
  \vcenter{\hbox{\rule{0pt}{\MatrixWritingSize}\makebox[#1]{}}}}
\newcommand{\MatrixBox}{\AnswerBox{\MatrixBoxWidth}{\MatrixSize}{%
  $\left(\WritingArea{\MatrixWritingSize}\right)$}}
\newcommand{\VectorBox}{\AnswerBox{\VectorWidth}{\MatrixSize}{%
  $\left(\WritingArea{\VectorWritingWidth}\right)$}}
\newcommand{\AugmentedMatrixBox}{\AnswerBox{\AugmentedBoxWidth}{\MatrixSize}{%
  $\setlength{\arraycolsep}{3pt}\left(\begin{array}{@{}c|c@{}}
    \WritingArea{\MatrixWritingSize}&\WritingArea{\AugmentedColumnWidth}
  \end{array}\right)$}}
\newcommand{\EigenvalueBox}{\BlankBox{\EigenvalueWidth}{\ShortHeight}}
\newcommand{\SolutionBox}{\BlankBox{\SolutionWidth}{\SolutionHeight}}
\newcommand{\DerivativeBox}{\BlankBox{\DerivativeWidth}{\ShortHeight}}
\newcommand{\WideDerivativeBox}{\BlankBox{\WideDerivativeWidth}{\ShortHeight}}
\newcommand{\VectorSolutionBox}{\BlankBox{\VectorSolutionWidth}{\MatrixSize}}
\newcommand{\LinearSolutionBox}{\BlankBox{\LinearSolutionWidth}{\LinearSolutionHeight}}
\newcommand{\ParticularSolutionBox}{\BlankBox{\ParticularSolutionWidth}{\LongSolutionHeight}}
\newcommand{\GeneralSolutionBox}{\BlankBox{\GeneralSolutionWidth}{\LongSolutionHeight}}

\begin{document}
\WorksheetHeading{A Forced $3\times3$ IVP}
Solve the IVP and determine behavior as $t\to\infty$. 
\[
  \vect{y}'=A\vect{y}+e^{-2t}\begin{pmatrix}-4\\-4\\-4\end{pmatrix}
  +t\begin{pmatrix}2\\2\\0\end{pmatrix}
  \qquad A=\begin{pmatrix}-4&2&4\\-2&1&3\\-4&3&3\end{pmatrix}
  \qquad\vect{y}(0)=\vect{y}_0=\begin{pmatrix}3\\1\\4\end{pmatrix}
\]
\begingroup
\renewcommand{\labelenumi}{\textbf{\alph{enumi}.}}
\setlength{\leftmargini}{1.4em}
\begin{enumerate}
\setlength{\itemsep}{5pt}
\setlength{\parsep}{5pt}
\setlength{\topsep}{0pt}
\item\textbf{Eigenvalues and eigenvectors}\par
$\displaystyle\lambda_1=\EigenvalueBox\quad
\lambda_2=\EigenvalueBox\quad
\lambda_3=\EigenvalueBox$
\par
$\displaystyle V=[\vect{v}_1\ \vect{v}_2\ \vect{v}_3]=\MatrixBox$

\item\textbf{Real complementary solution}\quad
$\vect{y}_c(t)=c_1\vect{y}_1(t)+c_2\vect{y}_2(t)+c_3\vect{y}_3(t)$\par
Write down three linearly independent real solutions
\par
$\displaystyle\vect{y}_1(t)=\SolutionBox\hspace{\SolutionGap}
\vect{y}_2(t)=\SolutionBox\hspace{\SolutionGap}
\vect{y}_3(t)=\SolutionBox$

\item\textbf{Exponential forcing} \quad$\vect{f}_1(t)=e^{-2t}(-4,-4,-4)^T$

Try $\vect{y}_{p_1}(t)=\vect{a}e^{-2t}$\quad$\displaystyle\vect{y}_{p_1}'(t)=\WideDerivativeBox$
\par Substitute into $\vect{y}_{p_1}'=A\vect{y}_{p_1}+\vect{f}_1$ and cancel the exponential.\par
$\displaystyle\MatrixBox\vect{a}=\VectorBox\qquad\vect{y}_{p_1}(t)=\VectorSolutionBox$

\newpage
\item\textbf{Linear forcing} \quad
	$\vect{f}_2(t)=t(2,2,0)^T$

Try $\vect{y}_{p_2}(t)=\vect{b}t+\vect{d}$\quad$\displaystyle\vect{y}_{p_2}'(t)=\DerivativeBox$
\par Substitute into $\vect{y}_{p_2}'=A\vect{y}_{p_2}+\vect{f}_2$ and match the coefficient of $t$ and the constant term and solve

$\displaystyle\vect{y}_{p_2}(t)=\LinearSolutionBox$

\item\textbf{Particular and general solutions}\par
$\displaystyle\vect{y}_p(t)=\vect{y}_{p_1}(t)+\vect{y}_{p_2}(t)=\ParticularSolutionBox$
\par
$\displaystyle\vect{y}(t)=\vect{y}_c(t)+\vect{y}_p(t)=\GeneralSolutionBox$

\item\textbf{Initial condition in GS $\vect{y}(t)=\vect{y}_c(t)+\vect{y}_{p_1}(t)+\vect{y}_{p_2}(t) $}

Write down the augmented matrix you would row reduce to solve for the 
constants $\vect{c}$ to satisfy $\vect{y}(0)=xxx$.  
\par$\displaystyle\AugmentedMatrixBox$

\item\textbf{IVP solution and leading behavior}\par
Solve for $\vect{c}$ and fill in the asymptotic behavior. 
\end{enumerate}
\endgroup
\end{document}
